\subsection{Seed variance and surprisal readings}

The bridge identities also admit a disciplined probabilistic reading. Since
\[
\lambda_C^2=\phiz(1-\phiz),
\]
the Cabibbo scale is exactly the standard-deviation scale of a Bernoulli variable with success probability $p=\phiz$:
\[
\lambda_C=\sqrt{p(1-p)}.
\]
Equivalently, the Fisher information of the Bernoulli family is
\[
I(p)=\frac{1}{p(1-p)}=\frac{1}{\lambda_C^2},
\]
so the Cabibbo bridge measures the inverse square-root Fisher scale of the retained seed. This is the precise information-theoretic statement. The present paper does \emph{not} identify $\lambda_C$ itself with Fisher information or with a Bures distance.

Likewise, whenever the electromagnetic closure is written with
\[
\ln\!\bigl((\phiz)^{-1}\bigr),
\]
the exact probabilistic reading is the surprisal of the seed event. That term should not be called the Shannon entropy of the full Bernoulli distribution, which would instead be
\[
H(\phiz)
=
-\phiz\ln\phiz-(1-\phiz)\ln(1-\phiz).
\]
Finally, the neutrino bridge relation
\[
\sin^2\theta_{13}=\phiz e^{-\gamma}
\]
is naturally described as Boltzmann-like weighting by the carrier norm $\gamma$, but the present draft does not define a full thermodynamic partition function at that point. These readings are therefore organizational interpretations of already established formulas, not new closure theorems.

\subsection{Exact parity-code and information-compression reading}

The strongest information-theoretic statement is no longer a loose analogy. By \Cref{prop:parity-code-family}, the one-family packet is exactly the even-parity subspace on five occupation bits,
\[
S^+
\cong
\mathrm{span}\Bigl\{|x\rangle:\sum_{i=1}^5 x_i\equiv 0 \pmod 2\Bigr\},
\]
so the carrier packet is kinematically the classical even-parity $[5,4,2]$ code. The split formula
\[
S^+
=
\bigl(\Lambda^{\mathrm{even}}E_3\otimes\Lambda^{\mathrm{even}}E_2\bigr)
\oplus
\bigl(\Lambda^{\mathrm{odd}}E_3\otimes\Lambda^{\mathrm{odd}}E_2\bigr)
\]
shows that the coarse parity bits of the $3$-sector and $2$-sector are exactly locked. For the maximally mixed state on the $16$ basis words, the two parity observables are perfectly correlated and therefore carry exactly one bit of mutual information.

The split weight enumerator in \Cref{prop:split-weight-enumerator} is the corresponding code enumerator with hypercharge weights. It packages the six admissible occupancy sectors into the single polynomial
\[
W(u,v)
=
1+3u^2+6uv+v^2+2u^3v+3u^2v^2.
\]
The hypercharge-compression result of \Cref{prop:family-hypercharge-compression} sharpens the same picture further. On the uniform distribution over the $16$ carrier basis states, the six $Y$-sectors have probabilities
\[
\frac1{16}(1,3,6,1,2,3),
\]
so the hypercharge entropy is
\[
H(Y)
=
-\sum_i p_i\log_2 p_i
=
\frac72-\frac34\log_2 3
\approx
\num{2.311}\ \text{bits}.
\]
Since the total carrier word takes $4$ bits to specify, the unresolved nonabelian degeneracy is
\[
4-H(Y)\approx \num{1.689}\ \text{bits}.
\]
Thus hypercharge is an efficient family-sector compressor, but not a lossless label for the individual basis word.

\begin{remark}[Two code layers: kinematic and admissible]
The code reading should be split into two layers. First, $S^+$ is exactly the classical even-parity $[5,4,2]$ code on five occupation bits. Second, on the commuting admissible branch the protected physical sector is
\[
\mathcal H_{\mathrm{adm}}
:=
\operatorname{Ran}(\Padm)
=
\ker \Kadm
=
\ker \Delta_{\mathrm{adm}},
\]
where the last equality uses the Hodge description of \Cref{lem:hodge-admissibility-projector}. In that second layer the penalties $K_\Sigma,K_c,K_\Theta,K_g$ act as stabilizer-like sector constraints, and $\Padm$ is the Hodge projector onto the admissible physical subspace. This is the precise bridge to quantum-information language used in the manuscript. It is not yet a theorem of a self-correcting quantum memory, because no noise model, logical algebra, or lifetime analysis is derived here.
\end{remark}

\subsection{Cut-and-project model-set realization}

The appendix-level quasicrystal language is kept outside the hard core, but it can be stated
exactly as a compatible model-set construction.

\begin{theorem}[Regular cut-and-project model set for the \texorpdfstring{$3+2$}{3+2} carrier]
\label{thm:cut-project-model-set}
Let
\[
L=\mathbb Z^5.
\]
Choose a direct sum decomposition
\[
\mathbb R^5=V_{\mathrm{vis}}\oplus V_{\mathrm{int}},
\qquad
\dim V_{\mathrm{vis}}=3,
\qquad
\dim V_{\mathrm{int}}=2,
\]
such that the projection of $L$ to $V_{\mathrm{int}}$ is dense. Let
\[
\pi_{\mathrm{vis}}:\mathbb R^5\to V_{\mathrm{vis}},
\qquad
\pi_{\mathrm{int}}:\mathbb R^5\to V_{\mathrm{int}}
\]
be the canonical projections, and let $W\subset V_{\mathrm{int}}$ be a compact regular window with
nonempty interior and boundary of measure zero. Define
\[
\Lambda_W
:=
\left\{
\pi_{\mathrm{vis}}(n)\;:\;
n\in L,\;
\pi_{\mathrm{int}}(n)\in W,\;
\sum_{i=1}^5 n_i\equiv 0\!\!\!\pmod 2
\right\}.
\]
Then $\Lambda_W$ is a regular model set. The parity constraint realizes the even-packet selection of
$S^+$, and the dimension split $3+2$ realizes the visible/internal carrier architecture.
\end{theorem}

\begin{proof}
This is the standard cut-and-project theorem for regular model sets. Density of the internal
projection and regularity of the window imply that $\Lambda_W$ is Delone and has pure point
diffraction. The congruence condition
\[
\sum_{i=1}^5 n_i\equiv 0\!\!\!\pmod 2
\]
selects the even parity class, matching the even exterior-packet sector of the carrier.
\end{proof}

\begin{remark}[Status of the quasicrystal reading]
The model-set theorem makes the quasicrystal appendix mathematically sharp, but the irrational split
and window are not derived from the hard TFPT core. The construction is therefore exact yet
deliberately kept appendix level.
\end{remark}

\section{Record algebra and pointer dynamics}
\label{app:record}

This appendix elaborates the main-text record constructions and the stable pointer-algebra package
used in the downstream closure discussion. Horizon thermality and prediction / transient semantics
are separated into the following appendix sections so that record theory, black-hole continuation,
and phenomenological event channels do not blur into one mixed extrapolation block.

\subsection{Minimal admissible vacuum}

At finite volume or with a declared cutoff, an admissible vacuum candidate is written as
\[
\rho_B
:=
\argmin_{\rho\in\Vadm}\tr(\rho\Hrel),
\]
with lower boundedness and coercivity understood exactly as in the main-text closure axioms. The stronger generative theorem path described in the conclusion would refine this object to the global admissible vacuum functional $\mathfrak V[g,\chigeo,\Phi,\nabla_F,\delta,\{L_f\}]$ whose variations also select the transport branch. In this appendix, however, the same vacuum object is used only as the background state for the record-algebra elaborations; no stronger cosmological claim is added here beyond the closed existence statement already declared in the main theorem stack.

\subsection{Record growth principle}

One useful coarse-grained record algebra is
\[
\Arec(t):=\bigoplus_{\Delta_i>\Hhub(t)}\mathcal A_i,
\qquad
S_{\mathrm{rec}}(t)=\log\dim\Arec(t).
\]
As $\Hhub(t)$ decreases, more stable channels open. The intended arrow-of-time claim is therefore
\[
\frac{dS_{\mathrm{rec}}}{dt}>0,
\]
but only as a record-growth principle, not as a replacement for detailed nonequilibrium dynamics.

\subsection{Observer algebra and Born-type evaluation}

An observer is modeled as a stable pointer subalgebra
\[
\Aobs\subset\Arec.
\]

\begin{theorem}[Born evaluation on the stable pointer algebra]
\label{thm:born-pointer-algebra}
Let $\{P_i\}\subset \Aobs$ be a complete family of orthogonal projectors. Then for every state
$\rho$ the unique additive, positive, and phase-invariant evaluation rule on this projector family is
\[
p_i=\tr(\rho P_i).
\]
\end{theorem}

\begin{proof}
In finite-dimensional sectors this is the standard Born rule characterization. The extra TFPT input
is not a different probability law but the selection of the stable pointer algebra by the gapped
admissible sector.
\end{proof}

With these projectors one may write the measurement channel
\[
\mathcal M_O(\rho)=\sum_i(P_i\rho P_i)\otimes |i\rangle\langle i|.
\]
The theorem is therefore a constrained appendix-level evaluation statement on admissible sectors,
not a new claim about measurement in arbitrary settings.

\section{Horizons, modular flow, and black hole thermality}
\label{app:horizons}

This appendix section records the horizon-side analytic input that supports the downstream modular
thermality program. In the present version the split-inclusion theorem is the firm appendix-level
statement, while the full modular KMS upgrade and compact-object thermality package remain visible
downstream closure targets.

\subsection{Split inclusion horizon approximation}

\begin{theorem}[Horizon approximation by split inclusion]
\label{thm:horizon-split-inclusion}
Let $\mathfrak A_{\mathrm{adm}}$ be the local admissible net reconstructed in
\Cref{thm:admissible-local-minkowski-net}, and assume in addition the split property. Then for
every nested triple
\[
\mathcal O_{\mathrm{out}}\Subset \mathcal O_{\mathrm{collar}}\Subset \mathcal O_{\mathrm{ext}}
\]
there exists a type-I factor $N$ such that
\[
\mathfrak A_{\mathrm{adm}}(\mathcal O_{\mathrm{out}})
\subset
N
\subset
\mathfrak A_{\mathrm{adm}}(\mathcal O_{\mathrm{ext}}).
\]
Hence on a suitable code subspace one obtains an approximate tensor factorization
\[
\mathcal H_{\mathrm{code}}
\cong
\mathcal H_{\mathrm{out}}\otimes\mathcal H_{\mathrm{hor}}\otimes\mathcal H_{\mathrm{in}}.
\]
\end{theorem}

\begin{proof}
This is the standard consequence of the split property for local quantum field theory nets. The
result supplies an algebraic horizon approximation without asserting an exact fundamental
factorization.
\end{proof}

If a horizon is modeled as an inner seam, the manuscript therefore uses only this split-based
approximation language rather than an exact microscopic tensor decomposition.

\section{Prediction semantics and transient event channels}
\label{app:prediction-semantics}

This appendix section collects the optional semantic decorations of the prediction map and the
status note for transient event channels. These items organize downstream interpretation, but they
do not add new upstream theorem arrows to the rigid core.

\subsection{Prediction semantics}

The main-text prediction map can be decorated by explicit experimental metadata as
\[
Q=(\rho_0,\mathcal O_Q,\Delta_Q,t,\mathcal C_Q,\mathcal R_Q),
\qquad
\Pred(Q)=\mathcal R_Q\!\left(\tr(\rho_tE_Q(\Delta)),\mathcal C_Q\right),
\]
with
\[
\rho_t=U(t)\rho_0U^\dagger(t).
\]
This is kept only as an optional formal extension.

\subsection{Transient event status}

The topological burst channel of \Cref{conj:topological_phase_slip_burst_channel} is kept at
programmatic-target status. What is fixed in the present version is only the admissible phase-slip
channel and its scaling form, not yet a complete defect dynamics or radiative-transfer derivation.

\section{Closed theoretical outputs and empirical comparison surface}
\label{app:global-particle-ledger}

This appendix collects the full Standard Model output of the boundary-polarized manuscript into a
theoretical surface and an empirical comparison surface. The purpose is not to flatten the
hierarchy but to keep theorem-level outputs, derived consequences, physical observables, scheme projections, and appendix continuations
visibly distinct:

\begin{tcolorbox}[colback=tfptblue!2, colframe=tfptblue!50, breakable,
    title={Status legend for the particle ledger}, colbacktitle=tfptblue!70]
\begin{itemize}[leftmargin=1.5em, topsep=3pt, itemsep=2pt]
\item \textbf{Theorem:} a theorem-level output on the closed branch.
\item \textbf{Derived consequence:} a structural output obtained by composing earlier theorems without adding a new scheme convention.
\item \textbf{Physical observable:} a branch-fixed observable read internally before any scheme choice.
\item \textbf{Scheme projection:} a row reported only after declared schemes, thresholds, or matching conventions.
\item \textbf{Appendix continuation:} appendix-level continuation not needed for the closed core.
\item \textbf{Programmatic target:} a precise downstream closure target recorded without a claimed proof in the present version.
\end{itemize}
Rows on the empirical comparison surface are compared only under the declared external
conventions once the physical observable layer has already been fixed.
\end{tcolorbox}

\begin{remark}[Counting convention used here]
\label{rem:counting-convention}
One chiral family from $S^+=\Lambda^{\mathrm{even}}E$ contains $16$ states:
$\nu^c$\,(1), $u^c$\,(3), $Q_L$\,(6), $e^c$\,(1), $L_L$\,(2), $d^c$\,(3).
With $\nu^c$ explicitly in the carrier family packet, each family has $6$ Weyl field types.
Three families give $\Oadm=3\times 16=48$ admissible chiral states.
The gauge sector has $12$ bosons ($\gamma$, $8$ gluons, $W^\pm$, $Z$), and the Higgs sector has
$1$ physical scalar. The total field-type count is therefore $3\times 6+4+1=23$ distinct types.
\end{remark}

\subsection{Table 1: Closed theoretical outputs}

\begin{table}[H]
\centering
\tiny
\setlength{\tabcolsep}{1.5pt}
\renewcommand{\arraystretch}{1.0}
\begin{tabular}{@{}>{\raggedright\arraybackslash}p{0.09\linewidth}>{\raggedright\arraybackslash}p{0.11\linewidth}>{\raggedright\arraybackslash}p{0.10\linewidth}>{\centering\arraybackslash}p{0.04\linewidth}>{\raggedright\arraybackslash}p{0.18\linewidth}>{\raggedright\arraybackslash}p{0.20\linewidth}>{\raggedright\arraybackslash}p{0.08\linewidth}@{}}
\toprule
Particle & Repr. & Ext.\ source & $Y$ & TFPT origin & Internal source & Status \\
\midrule
\multicolumn{7}{@{}l}{\textit{Fermion sector (one family $\times\,3$):}} \\
$\nu^c$ & $(1,1)_0$ & $\Lambda^0E$ & $0$ & carrier scaffold & hard carrier, $S^+$ & Theorem \\
$e^c$ & $(1,1)_1$ & $\Lambda^2E_2$ & $1$ & carrier scaffold & hard carrier, $S^+$ & Theorem \\
$u^c$ & $(\bar 3,1)_{-2/3}$ & $\Lambda^2E_3$ & $-\frac23$ & carrier scaffold & hard carrier, $S^+$ & Theorem \\
$d^c$ & $(\bar 3,1)_{1/3}$ & \makecell[l]{$\Lambda^2E_3$\\$\otimes\Lambda^2E_2$} & $\frac13$ & carrier scaffold & hard carrier, $S^+$ & Theorem \\
$Q_L$ & $(3,2)_{1/6}$ & $E_3\otimes E_2$ & $\frac16$ & carrier scaffold & hard carrier, $S^+$ & Theorem \\
$L_L$ & $(1,2)_{-1/2}$ & \makecell[l]{$\Lambda^3E_3$\\$\otimes E_2$} & $-\frac12$ & carrier scaffold & hard carrier, $S^+$ & Theorem \\
\midrule
\multicolumn{7}{@{}l}{\textit{Gauge boson sector:}} \\
$\gamma$ & $(1,1)_0$ & unbroken $U(1)_{\mathrm{em}}$ & $0$ & exact gauge direction & carrier algebra & Theorem \\
$g$ ($\times 8$) & $(8,1)_0$ & unbroken $SU(3)_c$ & $0$ & exact gauge direction & carrier algebra & Theorem \\
$W^\pm$ & $(1,3)_{\pm1}$ & SSB from $\Phi$ & $\pm1$ & bosonic index + Higgs mechanism & Dirac-anchor closure & Physical observable \\
$Z$ & $(1,1)_0$ & SSB from $\Phi$ & $0$ & bosonic index + Higgs mechanism & Dirac-anchor closure & Physical observable \\
\midrule
\multicolumn{7}{@{}l}{\textit{Scalar sector:}} \\
$H$ & $(1,2)_{1/2}$ & seam-even $\Phi$ & $\frac12$ & compact bosonic index & $\Ind(\mathcal B_{E_2}^+)=1$ & Theorem \\
\midrule
\multicolumn{7}{@{}l}{\textit{Structural outputs:}} \\
$N_{\mathrm{fam}}=3$ & $\mathcal H_F\cong F\cong\mathbb C^3$ & $L^2$-APS harmonic family modes & --- & harmonic family-mode theorem & punctured family section & Theorem \\
$\Oadm=48$ & $S^+\otimes \mathcal H_F$ & occupancy & --- & bridge corollary & $3\times 16$ & Theorem \\
$G_{\mathrm{car}}$ & \makecell[l]{$S(U(3)\!\times\!U(2))$\\$\cong[SU(3)\!\times\!SU(2)$\\$\times U(1)_Y]/\mathbb Z_6$} & carrier split & --- & internal carrier stabilizer theorem (\Cref{thm:global-sm-theorem}) & $E_3\oplus E_2$ & Theorem \\
local algebra & $\mathfrak{su}(3)\oplus\mathfrak{su}(2)\oplus\mathfrak u(1)$ & local Lie algebra & --- & differential of the internal carrier theorem & local sections of $E_3\oplus E_2$ & Derived consequence \\
\bottomrule
\end{tabular}
\caption{Standard Model content and origin map on the closed branch. The status column separates theorem outputs from physical observables and appendix-only scheme projections.}
\label{tab:sm-content-origin}
\end{table}

